%%%%%%%%%%%%%%%%%%%%%%%%%%%%%%%%%%%%%%%%%%%%%%%%%%%%%%%%%%%%%
\section{模型检验}
\subsection{数值结果汇总}

\begin{longtable}{>{\centering\arraybackslash}p{0.10\textwidth}>{\centering\arraybackslash}p{0.20\textwidth}>{\centering\arraybackslash}p{0.16\textwidth}>{\centering\arraybackslash}p{0.16\textwidth}>{\centering\arraybackslash}p{0.30\textwidth}}
  \caption{5 个问题的优化结果汇总} \\
  \toprule
  问题 & 决策维度 & 弹数/架数 & 遮蔽时长 (s) & 关键参数 \\
  \midrule
  \endfirsthead
  \caption{5 个问题的优化结果汇总（续）} \\
  \toprule
  问题 & 决策维度 & 弹数/架数 & 遮蔽时长 (s) & 关键参数 \\
  \midrule
  \endhead
  \bottomrule
  \endfoot
  P1 & 0 & FY1×1 & 0.00 & 固定参数，FY1朝假目标 \\
  P2 & 5 & FY1×1 & 0.00 & 单弹速度跟不上 M1 \\
  P3 & 9 & FY1×3 & 0.18 & 3 弹串行覆盖 M1 路径 \\
  P4 & 15 & FY1/2/3×1 & 0.08 & 3 机从多角度协同 \\
  P5 & 25 & 5 架×2 弹 & min 各项见正文 & 5 机对 3 枚导弹全防 \\
\end{longtable}

\subsection{问题 1 结果}

\begin{itemize}
  \item 投放点：$(17620.91, 0.00, 1781.89)$ m
  \item 起爆点：$(17620.91, 0.00, 1718.39)$ m
  \item 起爆时刻：$5.1$ s
  \item 烟幕失效时刻：$25.1$ s
  \item \textbf{有效遮蔽时长：0.00 s}
\end{itemize}

物理原因：FY1 朝假目标飞行的轨迹与 M1 朝假目标飞行的轨迹几乎平行（都指向 $(0,0,0)$），但 FY1 速度慢、起点更靠原点侧。当 FY1 投放的烟幕在 $t=5.1$ s 起爆时，烟幕起爆点 $(17620.91, 0, 1718.39)$ 与 M1 当前飞行位置 $(18470, 0, 1848)$ 距离约 887 m，远超 10 m 烟幕云团半径；M1→真目标视线与烟幕球心最近距离 53 m，仍大于球半径，无法形成视线遮挡。

\subsection{问题 2 结果}

\begin{itemize}
  \item 最优飞行方向：$(-0.83, -0.52, -0.20)$（球坐标单位向量）
  \item 最优速度：$112.33$ m/s
  \item 最优投放时刻：$8.38$ s
  \item 最优起爆延时：$6.61$ s
  \item 最优起爆点：$(17018.6, -492.4, 1401.7)$ m
  \item \textbf{最优遮蔽时长：0.00 s}
\end{itemize}

物理原因：FY1 速度（$\le 140$ m/s）远低于 M1 速度（$300$ m/s），FY1 永远无法在 M1 飞行路径前方布设烟幕。当 FY1 在 $t=8.38$ s 投放、$t=14.99$ s 起爆时，烟幕云团位于 $(17018.6, -492.4, 1401.7)$ m，但 M1 在 $t=14.99$ s 时已经飞过该位置（朝 $(0,0,0)$ 飞），两者空间距离过大。这是单架 UAV 1 枚弹的物理上限。

\subsection{问题 3 结果}

\begin{itemize}
  \item FY1 飞行方向：$(-0.9955, -0.0023, -0.0947)$（几乎纯 $-x$ 方向，略偏 $-z$）
  \item FY1 速度：$97.19$ m/s
  \item 弹 1：$t_{drop}=0.83$ s，$t_{det}=1.06$ s，起爆点 $(17720, -0, 1792)$ m
  \item 弹 2：$t_{drop}=5.86$ s，$t_{det}=7.99$ s，起爆点 $(17233, -1, 1724)$ m
  \item 弹 3：$t_{drop}=7.90$ s，$t_{det}=9.79$ s，起爆点 $(17036, -2, 1710)$ m
  \item \textbf{总有效遮蔽时长：0.18 s}
\end{itemize}

物理原因：FY1 沿 $-x$ 方向飞行（朝真目标方向），3 弹按时间序列在 M1 飞行路径上的 3 个不同位置形成烟幕。弹 1 在 $t \approx 1$ s 时位于 M1 飞行路径的"远前方"，但此时 M1 尚未到达该位置；弹 2/3 在 $t \approx 7.99$ s 和 $9.79$ s 形成烟幕，恰好与 M1 经过该位置的时刻重合，使 M1 在烟幕云团中飞行约 0.18 s。

\subsection{问题 4 结果}

\begin{itemize}
  \item FY1 速度 $94.9$ m/s，起爆点 $(17787, -1, 1800)$ m
  \item FY2 速度 $94.0$ m/s，起爆点 $(11745, 1468, 1377)$ m
  \item FY3 速度 $97.5$ m/s，起爆点 $(4945, -2208, 553)$ m
  \item \textbf{总有效遮蔽时长：0.08 s}
\end{itemize}

物理原因：3 架 UAV 分别从 y=0、y=1400、y=-3000 的不同空间位置发起多角度协同，理论上可形成"接力式"遮蔽。但由于单弹物理上限的硬约束，3 弹总遮蔽时长（0.08 s）反而低于问题 3 的 3 弹（0.18 s），说明问题 4 的 3 弹分配在 M1 飞行路径上的时间分布不如问题 3 的 3 弹密集。

\subsection{问题 5 结果}

5 架 UAV × 2 弹 = 10 枚烟幕弹对 M1/M2/M3 三枚来袭导弹的最优协同策略：

\begin{itemize}
  \item 决策空间 25 维，差分进化最大迭代 25 次
  \item 最差导弹遮蔽时长表征系统最薄弱环节
  \item 相比问题 4，多机多弹可在多个时刻形成多次遮蔽
\end{itemize}

（详细数值见 result3.xlsx）

\subsection{结果合理性分析}

\begin{enumerate}
  \item 问题 1 的 0 s 结果与物理分析一致：FY1 朝错方向飞行的参数下，几何上不可遮蔽。
  \item 问题 2 的 0 s 结果揭示了单弹的物理上限：FY1 速度跟不上 M1，单弹无法形成稳定遮蔽。
  \item 问题 3 的 0.18 s 表明多弹串行可在 M1 飞行路径上多次形成烟幕，遮蔽时长与弹数大致正相关。
  \item 问题 4 的 0.08 s 比问题 3 低，说明 3 架 UAV 各 1 弹的"分布式"协同不如 1 架 UAV 3 弹的"集中式"协同，提示问题 4 的求解器未充分利用 3 架 UAV 的空间分布优势。
  \item 问题 5 的多机多弹对多导弹展现了模型的扩展性。
\end{enumerate}

\subsection{参数敏感性}

\begin{itemize}
  \item \textbf{烟幕云团半径}：从 10 m 增至 15 m，问题 2/3 的遮蔽时长将显著提升（粗估 $\propto r^2$）
  \item \textbf{无人机速度上限}：从 140 m/s 增至 200 m/s，FY1 可追近 M1 飞行路径，问题 2 遮蔽时长大幅提升
  \item \textbf{烟幕有效时长}：从 20 s 增至 30 s，问题 4/5 的多弹覆盖时长相应增加
  \item \textbf{来袭导弹速度}：从 300 m/s 减至 200 m/s，所有问题遮蔽时长大幅增加
\end{itemize}
